\begin{thebibliography}{99}
\bibitem[1]{source1} Andreas Abel and Thierry Coquand. Failure of normalization in impredicative type theory with proof-irrelevant propositional equality. Log. Methods Comput. Sci., 16(2), 2020. URL: https://doi.org/10.23638/LMCS-16(2:14)2020 .
\bibitem[2]{source2} Yohji Akama, Stefano Berardi, Susumu Hayashi, and Ulrich Kohlenbach. An arithmetical hierarchy of the law of excluded middle and related principles. In 19th IEEE Symposium on Logic in Computer Science (LICS 2004), 14-17 July 2004, Turku, Finland, Proceedings, pages 192-201. IEEE Computer Society, 2004. URL: https://doi.org/10.1109/LICS.2004.1319613 .
\bibitem[3]{source3} Franco Barbanera and Stefano Berardi. Proof-irrelevance out of excluded-middle and choice in the calculus of constructions. Journal of Functional Programming, 6(3):519-525, 1996. URL: https://doi.org/10.1017/S0956796800001829 .
\bibitem[4]{source4} Josef Berger, Hajime Ishihara, and Peter Schuster. The weak koenig lemma, brouwer’s fan theorem, de morgan’s law, and dependent choice. Reports Math. Log., 47:63-86, 2012. URL: https://rml.tcs.uj.edu.pl/rml-47/3-Schuster.pdf .
\bibitem[5]{source5} Sophie Bernard, Cyril Cohen, Assia Mahboubi, and Pierre-Yves Strub. Unsolvability of the quintic formalized in dependent type theory. In Liron Cohen and Cezary Kaliszyk, editors, 12th International Conference on Interactive Theorem Proving, ITP 2021, Rome, Italy (Virtual Conference), June 29 - July 1, 2021, volume 193 of LIPIcs, pages 8:1-8:18. Schloss Dagstuhl - Leibniz-Zentrum für Informatik, 2021. URL: https://doi.org/10.4230/LIPIcs.ITP.2021.8 .
\bibitem[6]{source6} Errett Bishop. Foundations of Constructive Analysis. Mcgraw-Hill, New York, NY, USA, 1967.
\bibitem[7]{source7} Luitzen Egbertus Jan Brouwer. De onbetrouwbaarheid der logische principes. Tijdschrift voor Wijsbegeerte, 2:152-158, 1908. (In Dutch).
\bibitem[8]{source8} Ulrik Buchholtz, Tom De Jong, and Egbert Rijke. Epimorphisms and acyclic types in univalent foundations. The Journal of Symbolic Logic, pages 1-36, February 2025. URL: https://doi.org/10.1017/jsl.2024.76 .
\bibitem[9]{source9} Liron Cohen, Yannick Forster, Dominik Kirst, Bruno da Rocha Paiva, and Vincent Rahli. Separating Markov’s Principles. In LICS, pages 28:1-28:14. ACM, 2024. URL: https://doi.org/10.1145/3661814.3662104 .
\bibitem[10]{source10} Paul J. Cohen. The independence of the axiom of choice, 1963.
\bibitem[11]{source11} Thierry Coquand. Metamathematical investigations of a calculus of constructions. Technical Report RR-1088, INRIA, 1989. URL: https://hal.inria.fr/inria-00075471 .
\bibitem[12]{source12} Thierry Coquand and Gérard P Huet. The calculus of constructions. Information and Computation, 76(2/3):95-120, 1988. URL: https://doi.org/10.1016/0890-5401(88)90005-3 .
\bibitem[13]{source13} Radu Diaconescu. Axiom of choice and complementation. Proceedings of the American Mathematical Society, 51(1):176, August 1975. URL: https://doi.org/10.2307/2039868 .
\bibitem[14]{source14} Martín Hötzel Escardó. The Cantor–Schröder–Bernstein theorem for ∞-groupoids. Journal of Homotopy and Related Structures, 16(3):363-366, 2021. URL: https://doi.org/10.1007/s40062-021-00284-6 .
\bibitem[15]{source15} Solomon Feferman. Constructive theories of functions and classes. In Maurice Boffa, Dirk van Dalen, and Kenneth Mcaloon, editors, Logic Colloquium '78, volume 97 of Studies in Logic and the Foundations of Mathematics, pages 159-224. Elsevier, 1979. URL: https://doi.org/10.1016/S0049-237X(08)71625-2 .
\bibitem[16]{source16} Thiago Felicissimo, Yann Leray, Loïc Pujet, Nicolas Tabareau, Éric Tanter, and Théo Winterhalter. Definitional Proof Irrelevance Made Accessible. to be published at LICS 2026, 2026. URL: https://hal.science/hal-05474391 .
\bibitem[17]{source17} Yannick Forster. Church’s thesis and related axioms in Coq’s type theory. In Christel Baier and Jean Goubault-Larrecq, editors, 29th EACSL Annual Conference on Computer Science Logic, CSL 2021, January 25-28, 2021, Ljubljana, Slovenia (Virtual Conference), volume 183 of LIPIcs, pages 21:1-21:19. Schloss Dagstuhl - Leibniz-Zentrum für Informatik, 2021. URL: https://doi.org/10.4230/LIPIcs.CSL.2021.21 .
\bibitem[18]{source18} Yannick Forster and Felix Jahn. Constructive and synthetic reducibility degrees: Post’s problem for many-one and truth-table reducibility in Coq. In Bartek Klin and Elaine Pimentel, editors, 31st EACSL Annual Conference on Computer Science Logic, CSL 2023, Warsaw, Poland, February 13-16, 2023, volume 252 of LIPIcs, pages 21:1-21:21. Schloss Dagstuhl - Leibniz-Zentrum für Informatik, 2023. URL: https://doi.org/10.4230/LIPIcs.CSL.2023.21 .
\bibitem[19]{source19} Yannick Forster, Felix Jahn, and Gert Smolka. A Computational Cantor-Bernstein and Myhill’s Isomorphism Theorem in Constructive Type Theory (Proof Pearl). In Proceedings of the 12th ACM SIGPLAN International Conference on Certified Programs and Proofs, CPP ’23, pages 159-166. ACM, January 2023. URL: https://doi.org/10.1145/3573105.3575690 .
\bibitem[20]{source20} Yannick Forster, Dominik Kirst, and Niklas Mück. The Kleene-Post and Post’s Theorem in the Calculus of Inductive Constructions. In Aniello Murano and Alexandra Silva, editors, 32nd EACSL Annual Conference on Computer Science Logic, CSL 2024, Naples, Italy, February 19-23, 2024, volume 288 of LIPIcs, pages 29:1-29:20. Schloss Dagstuhl - Leibniz-Zentrum für Informatik, 2024. URL: https://doi.org/10.4230/LIPIcs.CSL.2024.29 .
\bibitem[21]{source21} Yannick Forster, Dominik Kirst, and Gert Smolka. On synthetic undecidability in Coq, with an application to the Entscheidungsproblem. In Proceedings of the 8th ACM SIGPLAN International Conference on Certified Programs and Proofs, CPP ’19, pages 38-51. ACM, January 2019. URL: https://doi.org/10.1145/3293880.3294091 .
\bibitem[22]{source22} Yannick Forster and Gert Smolka. Call-by-Value Lambda Calculus as a Model of Computation in Coq. Journal of Automated Reasoning, 63(2):393-413, October 2018. URL: https://doi.org/10.1007/s10817-018-9484-2 .
\bibitem[23]{source23} Makoto Fujiwara. Δ\ensuremath{{}^{0}}\ensuremath{{}_{1}} variants of the law of excluded middle and related principles. Archive for Mathematical Logic, 61(7–8):1113-1127, April 2022. URL: https://doi.org/10.1007/s00153-022-00827-5 .
\bibitem[24]{source24} Makoto Fujiwara and Ulrich Kohlenbach. Interrelation between weak fragments of double negation shift and related principles. The Journal of Symbolic Logic, 83(3):991-1012, September 2018. URL: https://doi.org/10.1017/jsl.2017.63 .
\bibitem[25]{source25} Gaëtan Gilbert. A type theory with definitional proof-irrelevance. (Une théorie des types avec insignifiance des preuves définitionnelle). PhD thesis, Mines ParisTech, France, 2019. URL: https://tel.archives-ouvertes.fr/tel-03236271 .
\bibitem[26]{source26} Jean-Yves Girard. Une extension de l'interprétation de gödel à l'analyse, et son application à l'élimination des coupures dans l'analyse et la théorie des types. In Proceedings of the Second Scandinavian Logic Symposium, volume 63 of Studies in Logic and the Foundations of Mathematics, pages 63-92. North-Holland, 1972.
\bibitem[27]{source27} Kurt Gödel. The consistency of the axiom of choice and of the generalized Continuum-Hypothesis. Proceedings of the National Academy of Sciences, 24(12):556-557, 1938. URL: https://doi.org/10.1073/pnas.24.12.556 .
\bibitem[28]{source28} N. Goodman and J. Myhill. Choice implies excluded middle. Mathematical Logic Quarterly, 24(25–30):461-461, January 1978. URL: https://doi.org/10.1002/malq.19780242514 .
\bibitem[29]{source29} Robert Harper and John C. Mitchell. Parametricity and variants of girard’s J operator. Inf. Process. Lett., 70(1):1-5, 1999. URL: https://doi.org/10.1016/S0020-0190(99)00036-8 .
\bibitem[30]{source30} Michael Hedberg. A coherence theorem for Martin-Löf’s type theory. Journal of Functional Programming, 8(4):413-436, 1998. URL: https://doi.org/10.1017/S0956796898003153 .
\bibitem[31]{source31} Arend Heyting. Die formalen Regeln der intuitionistischen Logik. Sitzungsberichte der Preußischen Akademie der Wissenschaften. Physikalisch-mathematische Klasse, pages 42-56, 1930.
\bibitem[32]{source32} Masasi Higasikawa. Partition principles and infinite sums of cardinal numbers. Notre Dame Journal of Formal Logic, 36(3):425-434, 1995. URL: https://doi.org/10.1305/NDJFL/1040149358 .
\bibitem[33]{source33} Martin Hofmann. Extensional Constructs in Intensional Type Theory. Springer London, 1997. URL: https://doi.org/10.1007/978-1-4471-0963-1 .
\bibitem[34]{source34} Paul Howard and Jean Rubin. Consequences of the Axiom of Choice. American Mathematical Society, June 1998. URL: https://doi.org/10.1090/surv/059 .
\bibitem[35]{source35} Antonius J. C. Hurkens. A simplification of Girard’s paradox, pages 266-278. Springer Berlin Heidelberg, 1995. URL: https://doi.org/10.1007/bfb0014058 .
\bibitem[36]{source36} Hajime Ishihara. Reverse mathematics in Bishop’s constructive mathematics. Philosophia Scientae, (CS 6):43-59, 2006. URL: https://doi.org/10.4000/philosophiascientiae.406 .
\bibitem[37]{source37} Dominik Kirst. Mechanised Metamathematics: An Investigation of First-Order Logic and Set Theory in Constructive Type Theory. PhD thesis, Saarland University, 2022. URL: https://doi.org/10.22028/D291-39150 .
\bibitem[38]{source38} Dominik Kirst and Haoyi Zeng. The Blurred Drinker Paradox: Constructive Reverse Mathematics of the Downward Löwenheim-Skolem Theorem. In 2025 40th Annual ACM/IEEE Symposium on Logic in Computer Science (LICS), pages 926-940. IEEE, June 2025. URL: https://doi.org/10.1109/lics65433.2025.00073 .
\bibitem[39]{source39} Stephen Cole Kleene. Recursive predicates and quantifiers. Transactions of the American Mathematical Society, 53(1):41, January 1943. URL: https://doi.org/10.2307/1990131 .
\bibitem[40]{source40} Stephen Cole Kleene and Richard Eugene Vesley. The foundations of intuitionistic mathematics: especially in relation to recursive functions, 1965.
\bibitem[41]{source41} Dominique Larchey-Wendling. Typing Total Recursive Functions in Coq, pages 371-388. Springer International Publishing, 2017. URL: https://doi.org/10.1007/978-3-319-66107-0\_24 .
\bibitem[42]{source42} Francis William Lawvere. Diagonal arguments and cartesian closed categories. In Lecture Notes in Mathematics, pages 134-145. Springer, 1969. URL: https://doi.org/10.1007/bfb0080769 .
\bibitem[43]{source43} Assia Mahboubi and Enrico Tassi. Mathematical Components. Zenodo, January 2021. URL: https://doi.org/10.5281/zenodo.4457887 .
\bibitem[44]{source44} Maria Emilia Maietti. About Effective Quotients in Constructive Type Theory, pages 166-178. Springer Berlin Heidelberg, 1999. URL: https://doi.org/10.1007/3-540-48167-2\_12 .
\bibitem[45]{source45} Maria Emilia Maietti. A minimalist two-level foundation for constructive mathematics. Annals of Pure and Applied Logic, 160(3):319-354, September 2009. URL: https://doi.org/10.1016/j.apal.2009.01.006 .
\bibitem[46]{source46} Maria Emilia Maietti and Giovanni Sambin. Toward a minimalistic foundation for constructive mathematics. In Laura Crosilla and Peter M. Schuster, editors, From sets and types to topology and analysis - Towards practicable foundations for constructive mathematics, volume 48 of Oxford logic guides. Oxford University Press, 2005.
\bibitem[47]{source47} Per Martin-Löf. Notes on Constructive Analysis. Almqvist \& Wiksell, Stockholm, 1968.
\bibitem[48]{source48} Per Martin-Löf. Intuitionistic type theory, volume 1 of Studies in proof theory. Bibliopolis, 1984.
\bibitem[49]{source49} Per Martin-Löf. 100 years of Zermelo’s axiom of choice: what was the problem with it? The Computer Journal, 49(3):345-350, December 2005. URL: https://doi.org/10.1093/comjnl/bxh162 .
\bibitem[50]{source50} Joan Rand Moschovakis. Disjunction, existence and \ensuremath{\lambda}-eliminability in formalized intuitionistic analysis. PhD thesis, University of Wisconsin-Madison, Madison, WI, USA, 1965.
\bibitem[51]{source51} Andrzej Mostowski. On definable sets of positive integers. Fundamenta Mathematicae, 34:81-112, 1947. URL: https://doi.org/10.4064/fm-34-1-81-112 .
\bibitem[52]{source52} John Myhill. Constructive set theory. Journal of Symbolic Logic, 40(3):347-382, September 1975. URL: https://doi.org/10.2307/2272159 .
\bibitem[53]{source53} Christine Paulin-Mohring. Inductive definitions in the system Coq rules and properties. In International Conference on Typed Lambda Calculi and Applications, pages 328-345. Springer, 1993. URL: https://doi.org/10.1007/BFb0037116 .
\bibitem[54]{source54} Josselin Poiret, Gaëtan Gilbert, Kenji Maillard, Pierre-Marie Pédrot, Matthieu Sozeau, Nicolas Tabareau, and Éric Tanter. All your base are belong to us: Sort polymorphism for proof assistants. Proc. ACM Program. Lang., 9(POPL):2253-2281, 2025. URL: https://doi.org/10.1145/3704912 .
\bibitem[55]{source55} Cécilia Pradic and Chad E. Brown. Cantor-Bernstein implies Excluded Middle. CoRR, abs/1904.09193, 2019. URL: https://arxiv.org/abs/1904.09193 .
\bibitem[56]{source56} Cécilia Pradic. The myhill isomorphism theorem does not generalize much, 2025. URL: https://doi.org/10.48550/arXiv.2507.05028 .
\bibitem[57]{source57} Loïc Pujet and Nicolas Tabareau. Observational equality: now for good. Proc. ACM Program. Lang., 6(POPL):1-27, 2022. URL: https://doi.org/10.1145/3498693 .
\bibitem[58]{source58} Loïc Pujet and Nicolas Tabareau. Impredicative Observational Equality. Proc. ACM Program. Lang., 7(POPL):2171-2196, 2023. URL: https://doi.org/10.1145/3571739 .
\bibitem[59]{source59} Fred Richman. Constructive Mathematics without Choice. In Peter Schuster, Ulrich Berger, and Horst Osswald, editors, Reuniting the Antipodes - Constructive and Nonstandard Views of the Continuum, pages 199-205. Springer Netherlands, Dordrecht, 2001. URL: https://doi.org/10.1007/978-94-015-9757-9\_17 .
\bibitem[60]{source60} The Rocq Development Team. The Rocq Prover, September 2025. URL: https://doi.org/10.5281/zenodo.17473943 .
\bibitem[61]{source61} Anne Sjerp Troelstra. Aspects of constructive mathematics. In Handbook of Mathematical Logic, volume 90 of Studies in Logic and the Foundations of Mathematics, pages 973-1052. North-Holland, Amsterdam, 1978.
\bibitem[62]{source62} Anne Sjerp Troelstra and Dirk van Dalen. Constructivism in mathematics. Vol. I. Studies in Logic and the Foundations of Mathematics, 26, 1988.
\bibitem[63]{source63} The Univalent Foundations Program. Homotopy Type Theory: Univalent Foundations of Mathematics. https://homotopytypetheory.org/book , Institute for Advanced Study, 2013.
\bibitem[64]{source64} Garyfallia Vafeiadou. A comparison of minimal systems for constructive analysis, 2018. URL: https://arxiv.org/abs/1808.00383 .
\bibitem[65]{source65} Mark van Atten. The Creating Subject, the Brouwer–Kripke Schema, and infinite proofs. Indagationes Mathematicae, 29(6):1565-1636, December 2018. URL: https://doi.org/10.1016/j.indag.2018.06.005 .
\bibitem[66]{source66} Wim Veldman. Brouwer’s Fan Theorem as an axiom and as a contrast to Kleene’s alternative. Archive for Mathematical Logic, 53(5–6):621-693, June 2014. URL: https://doi.org/10.1007/s00153-014-0384-9 .
\bibitem[67]{source67} Louis Warren, Hannes Diener, and Maarten McKubre-Jordens. The drinker paradox and its dual. CoRR, abs/1805.06216, 2018. URL: https://arxiv.org/abs/1805.06216 .
\bibitem[68]{source68} Benjamin Werner. On the strength of proof-irrelevant type theories. Logical Methods in Computer Science, Volume 4, Issue 3, September 2008. URL: https://doi.org/10.2168/lmcs-4(3:13)2008 .
\bibitem[69]{source69} Ernst Zermelo. Beweis, dass jede Menge wohlgeordnet werden kann. Mathematische Annalen, 59(4):514-516, 1904. URL: https://doi.org/10.1007/bf01445300 .
\bibitem[70]{source70} Ernst Zermelo. Untersuchungen über die grundlagen der mengenlehre. i. Mathematische Annalen, 65(2):261-281, 1908. URL: https://doi.org/10.1007/bf01449999 .
\end{thebibliography}
